% ============ Appendix to §13. Determinantal ideals (原书 pp.101-103) ============
\section*{Appendix to \S\ref{sec:13}. Determinantal ideals}
\phantomsection
\addcontentsline{toc}{section}{Appendix to \S\ref{sec:13}. Determinantal ideals}

Let $M = (a_{ij})$ be an $r \times s$ matrix $(r \le s)$ with elements $a_{ij}$ in a Noetherian ring
$A$, and let $I_t$ be the ideal of $A$ generated by the $t \times t$ minors (that is
$t$-determinants) of $M$. When $t = r$ and $A$ is a polynomial ring
$k[X_1, \ldots, X_s]$ over a field $k$, Macaulay proved that all the prime divisors of
$I_t$ have height $\le s - r + 1$ ([Mac], p. 54). In his Ph.D. thesis, Eagon
generalised this result as follows: for an arbitrary Noetherian ring $A$, every
minimal prime divisor of $I_t$ has height $\le (r - t + 1)(s - t + 1)$. The following
ingenious proof is taken from Eagon--Northcott [1]. We begin with some
preliminary observations.

The following operations on a matrix $M$ with elements in a ring $A$ are
called \emph{elementary row operations}: (1) permutation of the rows; (2) replacing
$C_i$ by $uC_i + vC_j$, where $C_i$ and $C_j$ $(i \neq j)$ are two distinct rows of $M$, $u$ is a unit
of $A$ and $v$ is an element of $A$; elementary column operations are defined
similarly. The ideal $I_t$ does not change under these operations. Now, if an
element of $M$ is a unit in $A$, we can transform $M$ by a finite number of
elementary row and column operations to the following form:
\[
  \begin{pmatrix}
    1 & 0 & \cdots & 0 \\
    0 & & & \\
    \vdots & & N & \\
    0 & & &
  \end{pmatrix},
\]
and $I_t$ is equal to the ideal of $A$ generated by the $(t - 1) \times (t - 1)$ minors of $N$.

\begin{lemma}\label{lem:13.1}
Let $(A, P)$ be a Noetherian local ring and set $B = A[X]$. Let $J$ be a
$P$-primary ideal of $A$ and $J'$ an ideal of $B$ such that $J' \subset PB$ and
$J' + XB = JB + XB$. Then $PB$ is a minimal prime divisor of $J'$.
\end{lemma}

\begin{proof}
$PB + XB$ is a maximal ideal of $B$, and is the radical of $JB + XB =
J' + XB$. Thus, in the ring $B/J'$ we have that $(PB + XB)/J'$ is a minimal
prime divisor of the principal ideal $(J' + XB)/J'$. Hence $\htop((PB + XB)/
J') = 1$. Since $PB \subset J'$ is a prime ideal in $B/J'$, we have $\htop(PB/J') = 0$.
\end{proof}

\begin{theorem}[Eagon]\label{thm:13.10}
Let $A$ be a Noetherian ring and $M$ be an $r \times s$
matrix $(r \le s)$ of elements of $A$. Let $I_t$ be the ideal of $A$ generated by the $t \times t$
minors of $M$. If $P$ is a minimal prime divisor of $I_t$, then we have
\[
  \htop P \le (r - t + 1)(s - t + 1).
\]
\end{theorem}

\begin{proof}
Induction on $r$. When $r = 1$ we have $t = 1$, and so $(r - t + 1)(s - t +
1) = s$. The ideal $I_t$ is generated by $s$ elements, so that the assertion is just
the principal ideal theorem (Theorem 5) in this case. Next assume that $r > 1$.
Localising at $P$ we may assume that $A$ is a local ring with maximal ideal $P$,
and that $I_t$ is $P$-primary.

If $t = 1$, then $I_t$ is generated by $rs$ elements and $(r - t + 1)(s - t + 1) = rs$,
so our assertion holds also for this case. Therefore we assume $t > 1$. If at
least one of the elements of $M$ is a unit of $A$, then by what we said above, $I_t$ is
generated by $(t - 1) \times (t - 1)$ minors of a $(r - 1) \times (s - 1)$ matrix, and again
we are done. Therefore we assume that all the elements of $M$ are in $P$. Now
comes the brilliant idea. Let $M'$ be the matrix with elements in $B = A[X]$
obtained from $M$ by replacing $a_{11}$ by $a_{11} + X$, and let $I'$ be the ideal of $B$
generated by the $t \times t$ minors of $M'$. Since $t > 1$ and $a_{ij} \in P$ for all $i$ and $j$ we
have $I' \subset PB$. We also have $I' + XB = I_tB + XB$ since both sides have the
same image in $B/XB = A$. Therefore $PB$ is a minimal prime divisor of $I'$ by
the lemma. Since the element $a_{11} + X$ of $M'$ is not in $PB$, we have
$\htop PB \le (r - t + 1)(s - t + 1)$ by our previous argument. Since $\htop PB = \htop P$,
as we can see by Theorems 4 and 5, we are done.
\end{proof}
